% Part:sets-functions-relations
% Chapter: size-of-sets
% Section: non-enumerability

\documentclass[../../../include/open-logic-section]{subfiles}

\begin{document}

\olfileid[ar]{sfr}{siz}{nen}

\olsection{مجموعات \printtoken{S}{nonenumerable}}

\begin{editorial}
  يثبت هذا القسم أن $\Bin^\omega$ و$\Pow{\PosInt}$ غير قابلتين
  للتعداد، باستعمال التعريف الوارد في \olref[enm]{sec}. وقد صُمم
  ليكون أبسط قليلًا وأكثر تفصيلًا قليلًا من الصيغة الواردة في
  \olref[enm-alt]{sec}.
\end{editorial}

بعض المجموعات، مثل مجموعة الأعداد الصحيحة الموجبة~$\PosInt$، لا
نهائية. وقد رأينا حتى الآن أمثلة لمجموعات لا نهائية كانت كلها
!!{enumerable}. لكن توجد أيضًا مجموعات لا نهائية لا تتمتع بهذه
الخاصية. وتسمى هذه المجموعات \emph{!!{nonenumerable}}.

ولعل وجود مجموعات !!{nonenumerable} مدعاة للدهشة أصلًا. فلكل مجموعة
!!{enumerable}~$A$ توجد دالة~!!a{surjective}
$f \colon \PosInt \to A$. أما إذا كانت مجموعة ما !!{nonenumerable}،
فلا توجد مثل هذه الدالة. أي لا يمكن لأي دالة ترسل !!{element}s~$\PosInt$،
وهي لا نهائية العدد، إلى~$A$ أن تستنفد~$A$ كلها. ومن ثم تكون
!!{element}s~$A$ «أكثر عددًا» من الأعداد الصحيحة الموجبة، على الرغم
من لا نهائيتها.

فكيف نثبت أن مجموعة ما !!{nonenumerable}؟ لا بد أن نبين أنه لا يمكن
وجود دالة شاملة من ذلك النوع. وبصورة مكافئة، لا بد أن نبين أن عناصر~$A$
لا يمكن تعدادها في قائمة لا نهائية أحادية الاتجاه. وأفضل طريقة لذلك
هي إثبات أن كل قائمة من !!{element}s~$A$ لا بد أن تغفل عنصرًا واحدًا
على الأقل؛ أو أن أي دالة $f\colon \PosInt \to A$ لا يمكن أن تكون
!!{surjective}. ويمكننا فعل ذلك باستعمال \emph{طريقة كانتور القطرية}.
فإذا أعطينا قائمة من !!{element}s~$A$، ولتكن $x_1$، $x_2$، \dots،
بنينا عنصرًا آخر من~$A$ يستحيل، بحكم طريقة بنائه، أن يكون في تلك
القائمة.

ومثالنا الأول هو المجموعة~$\Bin^\omega$، أي مجموعة جميع المتتاليات
اللانهائية الخالية من الفجوات والمؤلفة من~$0$ و~$1$.

\begin{thm}
\ollabel{thm:nonenum-bin-omega}
$\Bin^\omega$~!!{nonenumerable}.
\end{thm}

\begin{proof}
افترض، طلبًا للتناقض، أن $\Bin^\omega$ !!{enumerable}، أي افترض أن
ثمة قائمة $s_{1}$، $s_{2}$، $s_{3}$، $s_{4}$، \dots{} لجميع
!!{element}s~$\Bin^\omega$. وكل~$s_i$ منها متتالية لا نهائية
من~$0$ و~$1$. ولنرمز إلى العنصر ذي الرتبة~$j$ من المتتالية ذات
الرتبة~$i$ في هذه القائمة بـ$s_i(j)$. وعندئذ تكون المتتالية ذات الرتبة~$i$، أي
$s_i$، هي
\[
s_i(1), s_i(2), s_i(3), \dots
\]

يمكننا ترتيب هذه القائمة، وترتيب عناصر كل متتالية~$s_i$ فيها،
في مصفوفة:
\[
\begin{array}{c|c|c|c|c|c}
& 1 & 2 & 3 & 4 & \dots \\\hline
1 & \mathbf{s_{1}(1)} & s_{1}(2) & s_{1}(3) & s_1(4) & \dots \\\hline
2 & s_{2}(1)& \mathbf{s_{2}(2)} & s_2(3) & s_2(4) & \dots \\\hline
3 & s_{3}(1)& s_{3}(2) & \mathbf{s_3(3)} & s_3(4) & \dots \\\hline
4 & s_{4}(1)& s_{4}(2) & s_4(3) & \mathbf{s_4(4)} & \dots \\\hline
\vdots & \vdots & \vdots & \vdots & \vdots & \mathbf{\ddots}
\end{array}
\]
تعطي العناوين على الجانب رقم المتتالية في القائمة $s_1$، $s_2$،
\dots؛ أما الأعداد في الأعلى فتوسم !!{element}s المتتاليات المفردة.
فمثلًا، يرمز $s_{1}(1)$ إلى العدد، $0$ أو~$1$، أيًا يكن، الذي هو أول
!!{element} في المتتالية $s_{1}$، وهكذا.

والآن نبني متتالية لا نهائية، $\overline{s}$، من~$0$ و~$1$ يستحيل
أن تكون في هذه القائمة. وسيعتمد تعريف $\overline{s}$ على القائمة
$s_1$، $s_2$، \dots. فكل قائمة لا نهائية من متتاليات لا نهائية
من~$0$ و~$1$ تولد متتالية لا نهائية~$\overline{s}$ مضمونة الغياب عن
القائمة.

ولتعريف $\overline{s}$، نعين جميع !!{element}sها، أي نعين
$\overline{s}(n)$ لكل $n \in \PosInt$. ونفعل ذلك بقراءة قطر المصفوفة
أعلاه من أوله إلى آخره (ومن هنا اسم «الطريقة القطرية»)، ثم نغير كل
$1$ إلى~$0$، وكل $0$ إلى~$1$. وبصورة أكثر تجريدًا، نعرّف
$\overline{s}(n)$ ليكون $0$ أو $1$ بحسب ما إذا كان !!{element} ذو
الرتبة~$n$ على القطر، أي $s_n(n)$، هو $1$ أو $0$.
\[
\overline{s}(n) =
\begin{cases}
1 & \text{إذا كان $s_{n}(n) = 0$}\\
0 & \text{إذا كان $s_{n}(n) = 1$}.
\end{cases}
\]
وإذا كنت تفضل الصيغ على التعريف بالحالات، فيمكنك أيضًا تعريف
$\overline{s}(n) = 1 - s_n(n)$.

ومن الواضح أن $\overline{s}$ متتالية لا نهائية من~$0$ و~$1$، إذ ليست
إلا المتتالية المعاكسة لمتتالية قيم~$0$ و~$1$ الظاهرة على قطر
مصفوفتنا. ومن ثم فإن $\overline{s}$ !!a{element} من~$\Bin^\omega$.
لكنها لا يمكن أن تكون في القائمة $s_1$، $s_2$، \dots{} فلماذا؟

لا يمكن أن تكون أول متتالية في القائمة، أي $s_1$، لأنها تختلف عن
$s_1$ في أول !!{element}. فأيًا تكن قيمة $s_1(1)$، عرّفنا
$\overline{s}(1)$ لتكون نقيضها. ولا يمكن أن تكون ثاني متتالية في
القائمة، لأن $\overline{s}$ تختلف عن $s_2$ في العنصر الثاني: فإذا
كان $s_2(2)$ يساوي $0$، فإن $\overline{s}(2)$ يساوي $1$، والعكس
بالعكس. وهكذا.

وبصورة أدق: لو كانت $\overline{s}$ في القائمة، لوجد عدد~$k$ بحيث
$\overline{s} = s_{k}$. وتتطابق متتاليتان إذا وفقط إذا اتفقتا في كل
موضع، أي إنه، لأي~$n$، $\overline{s}(n) =
s_{k}(n)$. ومن ثم، على وجه الخصوص، إذا أخذنا $n = k$ حالةً خاصة، لزم
أن تكون $\overline{s}(k) = s_{k}(k)$. لكن $s_k(k)$ إما $0$ أو~$1$.
فإن كان $0$، وجب أن تكون $\overline{s}(k)$ مساويةً لـ~$1$---فهكذا
عرّفنا $\overline{s}$. أما إذا كان $s_k(k) = 1$، فعندئذ، وبسبب طريقة
تعريفنا $\overline{s}$ أيضًا، يكون $\overline{s}(k) = 0$. وفي كلتا
الحالتين $\overline{s}(k) \neq s_{k}(k)$.

بدأنا بافتراض وجود قائمة من !!{element}s~$\Bin^\omega$، هي $s_1$،
$s_2$، \dots{} وبنينا من هذه القائمة متتالية~$\overline{s}$ أثبتنا
أنها لا يمكن أن تكون فيها. لكنها بالتأكيد متتالية من~$0$ و~$1$ إذا
كانت جميع~$s_i$ متتاليات من~$0$ و~$1$، أي إن $\overline{s} \in
\Bin^\omega$. وهذا يبين، على وجه الخصوص، أنه لا يمكن وجود قائمة تضم
\emph{جميع} !!{element}s~$\Bin^\omega$، لأننا نستطيع أن نبني من أي
قائمة من ذلك النوع متتالية~$\overline{s}$ مضمونة الغياب عنها أيضًا؛
ومن ثم يؤدي افتراض وجود قائمة لجميع المتتاليات في~$\Bin^\omega$ إلى
تناقض.
\end{proof}

\begin{explain}
تسمى طريقة البرهان السابقة «الاستدلال القطري» لأنها تستعمل قطر
المصفوفة لتعريف~$\overline{s}$. ولا يلزم أن يتضمن الاستدلال القطري
وجود مصفوفة: إذ يمكننا إثبات أن بعض المجموعات ليست !!{enumerable} باستعمال
فكرة مماثلة حتى حين لا تكون ثمة مصفوفة ولا قطر فعلي.
\end{explain}

\begin{thm}
\ollabel{thm:nonenum-pownat}
$\Pow{\PosInt}$ ليست !!{enumerable}.
\end{thm}

\begin{proof}
نسير بالطريقة نفسها، فنثبت أنه لكل قائمة من المجموعات الجزئية
من~$\PosInt$ توجد مجموعة جزئية من $\PosInt$ لا يمكن أن تكون في القائمة.
افترض أن لدينا القائمة الآتية من المجموعات الجزئية من~$\PosInt$:
\[
Z_{1}, Z_{2}, Z_{3}, \dots
\]
نعرّف الآن مجموعة~$\overline{Z}$ بحيث إنه، لكل $n \in \PosInt$،
$n \in \overline{Z}$ إذا وفقط إذا كان $n \notin Z_{n}$:
\[
\overline{Z} = \Setabs{n \in \PosInt}{n \notin Z_n}
\]
ومن الواضح أن $\overline{Z}$ مجموعة من الأعداد الصحيحة الموجبة، لأن
كل~$Z_n$ كذلك بحسب الفرض؛ ومن ثم $\overline{Z} \in
\Pow{\PosInt}$. لكن $\overline{Z}$ لا يمكن أن تكون في القائمة. ولإثبات
ذلك، سنبين أنه، لكل $k \in \PosInt$، $\overline{Z} \neq
Z_k$.

ليكن إذن $k \in \PosInt$ اعتباطيًا. لقد عرّفنا $\overline{Z}$ بحيث
إنه، لأي $n \in \PosInt$، يكون $n \in \overline{Z}$ إذا وفقط إذا
كان $n \notin Z_n$. وعلى وجه الخصوص، إذا أخذنا $n=k$، كان
$k \in \overline{Z}$ إذا وفقط إذا كان $k \notin Z_k$. لكن هذا يبين
أن $\overline{Z} \neq Z_k$، لأن $k$ !!a{element} من إحداهما لا من
الأخرى؛ ولذلك تختلف $\overline{Z}$ و$Z_k$ في !!{element}sهما. ولما
كان $k$ اعتباطيًا، لم تكن $\overline{Z}$ في القائمة $Z_1$، $Z_2$،
\dots.
\end{proof}

\begin{explain}
لم يذكر البرهان السابق قطرًا، لكن يمكنك أن تتصوره متضمنًا قطرًا إذا
صورته على النحو الآتي: تصور المجموعات $Z_1$، $Z_2$، \dots، مكتوبةً
في مصفوفة، حيث يدرج كل !!{element}~$j \in Z_i$ في العمود ذي
الرتبة~$j$. ولتكن المجموعات الأربع الأولى في تلك القائمة
$\{1,2,3,\dots\}$، $\{2, 4, 6, \dots\}$، $\{1,2,5\}$،
و$\{3,4,5,\dots\}$. عندئذ تبدأ المصفوفة هكذا:
\[
\begin{array}{r@{}rrrrrrr}
  Z_1 = \{ & \mathbf{1}, & 2, & 3, & 4, & 5, & 6, & \dots\}\\
  Z_2 = \{ &  & \mathbf{2}, &  & 4, &  & 6, & \dots\}\\
  Z_3 = \{ & 1, & 2, &  &  & 5\phantom{,} &  & \}\\
  Z_4 = \{ &  &  & 3, & \mathbf{4}, & 5, & 6, & \dots\}\\
  \vdots & & & & & \ddots
\end{array}
\]
وعندئذ تكون $\overline{Z}$ المجموعة الناتجة من النزول على القطر، مع
حذف أي أعداد تظهر على القطر، وإدراج قيم~$j$ التي توجد عندها فجوة في
صف/عمود المصفوفة ذي الرتبة~$j$. وفي الحالة أعلاه نحذف $1$ و$2$،
وندرج~$3$، ونحذف~$4$، وهكذا.
\end{explain}

\begin{prob}
أثبت أن $\Pow{\Nat}$ !!{nonenumerable} بحجة قطرية.
\end{prob}

\begin{prob}\label{sfr:siz:nen:prob:f-posint}
أثبت، بحجة قطرية صريحة، أن مجموعة الدوال $f \colon \PosInt \to \PosInt$
!!{nonenumerable}. أي أثبت أنه إذا كانت $f_1$، $f_2$، \dots، قائمة
من الدوال، وكانت كل $f_i\colon
\PosInt \to \PosInt$، فثمة دالة $\overline{f}\colon \PosInt \to
\PosInt$ ليست في هذه القائمة.
\end{prob}

\end{document}
